\begin{workedexample}[Worked Example: EVM from SNR]
For a link degraded only by additive noise, the error-vector magnitude is
$\text{EVM}_{\text{rms}} \approx 10^{-\text{SNR}/20}$. At $\text{SNR} = 25\ \text{dB}$,
\[ \text{EVM}_{\text{rms}} \approx 10^{-25/20} = 0.056 = 5.6\%. \]
A $64$-QAM link typically needs EVM below about $2\%$ (SNR $>34\ \text{dB}$), so
this $25\ \text{dB}$ link would support only a lower-order constellation.
\end{workedexample}

\begin{keytakeaways}
\begin{itemize}[leftmargin=1.4em,itemsep=2pt]
\item I/Q carries two independent streams: $s = I\cos\omega_c t - Q\sin\omega_c t$.
\item Gain/phase imbalance $\to$ image; DC offset $\to$ carrier tone; LO leakage $\to$ carrier feedthrough.
\item EVM ties quality to SNR: $\text{EVM}\approx 10^{-\text{SNR}/20}$; it captures all impairments together.
\end{itemize}
\end{keytakeaways}

\begin{exercises}
\item What EVM corresponds to an SNR of $30\ \text{dB}$?
\item What impairment produces a spectral image of the wanted signal?
\item Why is EVM a more complete quality metric than SNR alone?
\end{exercises}

\furtherreading{Razavi~\cite{razavi} (ch.~4); Steer~\cite{steer}.}
